\chapter*{Introduction}
\addcontentsline{toc}{chapter}{Introduction}


for high-dimensional problems (e.g., $d > 5$), the exponential growth in storage
and computational complexity severely limits classical finite element and finite
difference methods


we first demonstrate the curse of dimensionality by using the finite differentce method on an ever-dimensionally-increasing pde system.

pde intro
showcase the curse
sdes
sparse grids



show on relatiely simple pde where sparse grids can sparse grids can be competitive


PDEs for modelling physical phenomena

curse of dimensionality

Machine learning in Scientific computing

Partial differential equations (PDEs) are essential mathematical 
tools for describing complex natural phenomena and engineering sys-tems
Traditional numerical meth-ods, such as the finite element method (FEM)
and the finite difference method (FDM), rely on spatial and temporal discretization to ap-proximate solutions by
constructing fine computational meshe
\cite{systematic_review}

increasingly used for modeling and simulating physical systems
despite some empirical success, PINNs are still facing many challenges
that define open areas for research and
further methodological advancements.
In recent years, there have been numerous studies focusing on improving the
performance of PINNs, mostly by
- neural network architectures
\cite{experts_guide}
fPINNS, XPINNS, varPINNS, ...


\cite{systematic_review}

In the past few years, solving PDEs using machine and deep learning methods has emerged as a new and rapidly growing field of scientific machine learning.
Leveraging the universal approximation theorem \cite{uni_approx} and high expressivity of neural networks, the 

first by Lagaris et. at. \cite{lagaris}

Karniadakis in 2019 changes the game \cite{comprehensive_analysis}

climate modelling etc.


data + physs loss
solve inverse problems - value of a fucntion or a coefficient from the data

In combine machine learning with partial differential equation is a unique and simple way that was intuitive to both the machine learning and pde solving communities.
https://www.sciencedirect.com/science/article/abs/pii/S0021999118307125

flexibility of the arch


We focues on semi large dims

very large dims possible with
- 100s 1000s but sol at a point:

for very high dims
- multilevel picard method
- deep backward storchastic eq?
- deep galerkin
on the machine learning side

and feynman-kac monte carlo formulation


no systematic or even mention of existing sparse grids approach

here: diffusin eq. - spread - fokker planck, etc.
commonly probabilistic setting.
high-dimensionality naturally arrises in FP

Sparse grids
- basically FD but on carfully chosed low discret points meshes
- by on bounded seconds ders able to make tradeoff tham make it possible to scale to higher dims


This thesis is structured as follows.
First, backgroung on diffusion pdes and stochasic differential equations
Issues we encounter when naively scaling to higher dimensions
Sparse grids
feed-forward neural networks
PINNs - specifics etc and simple experiments demostraining they scaling to higher dims
Results:
simple heat eq
more challenging smoluchowski
initial conditions uhlenbeck processes.



curse of dim - term coined by bellman in 1961

\chapter{Introduction}

%% PDEs and diffusion pdes
Partial differential equations are one of the core tools in mathematical physics,
modelling wide range of physical and engineering phenomena.
%including heat transfer, diffusion, fluid flow, wave propagation, biological systems, etc.
An important subclass form diffusion-type equations, describing the evolution of quantities
that spread over time, such as temperature, concentration,
probability density, or chemical particles in a medium.
Specific examples include the heat equation, the diffusion equation, and the Fokker-Planck equation.

%% standard mesh FEM etc
Standard approaches for solving such equations numerically
involve grid-based methods such as Finite difference,
Finite element, and Finite volume methods.
Due to their transparent construction
and theoretical convergence estimates,
those methods are a de facto standard.
However, their applicability is in practice limited to relatively
low-dimensional problems,
as scaling them up to higher dimensions 
leads to an exponentially rise
in the number of discretization points.
Thus quickly approaching the memory limitation of existing hardware.
%lead to the so called the curse of dimensionality,
%However, they suffer from the curse of dimensionality,

%% SG
Several approaches have been developed to combat this challenge.
One of which are so called sparse grids and notably,
the sparse grid combination technique.
Sparse grids combination technique circumvents the curse of dimensionality
by constructing a combinations of lower-dimensional anisotropic grids,
thus avoiding the full tensor-product construction of the standard approach.
Under suitable smoothness assumptions on the solution, they provide
a optimal compromise between approximation accuracy and computational cost,
making them attractive for moderately high-dimensional problems.
Sparse grids retain a strong connection to classical methods,
and the formulation of the combiantion technique
allows directly reusing existing solvers.
In addition, the special construction of the combination technique
allows it to be run in parallel.

%% Scientific machine learnaing
From the other directiion, scientific machine learning active research field.
The universal approximation theorem \cite{uni_approx} guarantees that,
under mild assumptions on the network architecture,
a sufficiently wide feed-forward network can
approximate any continuous function on a compact domain to arbitrary precision.
As such, neural networks
have beel applied on problems ranging from climate modelling, biological system, to fusion simulations.
\cite{hal_hands_on_pinns}
%Apart from their empirical success, recent work has
%given rise to many theoretical developments.

%% PINNS history
Early neural-network approaches to differential
equations were first proposed by Lagaris et al. \cite{lagaris} in 1999.
However, it was not until 2019, when
the introduction of physics-informed neural networks
by of Raissi, Perdikaris, and Karniadakis \cite{pinns_original}
caused a revolution in the field.

%% PINNS basic
One of the main advantages of PINNs is their simplicity and ease of 
integration into the existing neural network toolkit.
In comparison to the standard solvers mentioned above
PINNs are a mesh-free method.
Instead of constructing a full computational grid,
the points can be sampled randomly on the fly.
This gives PINNs a natural appeal for high-dimensional
problems, where traditional mesh-based methods suffer from exponentially
increasing memory requirements.

%% modern PINNS
Nevertheless, PINNs also have important limitations.
As is the case with standard neural networks,
their performance % mention runtime
significanlly depends on the choice of network architecture,
and hyperparameters.
Although many new model architercutal variants and extensions have been proposed,
including architectures such as fPINNs, XPINNs, Score PINN, and variational PINNs, 
just to name a few,
a reliable and efficient use of PINNs in higher-dimensional
settings still remains an active area of research.

%very large dims possible with
%- 100s 1000s but sol at a point:
%
%for very high dims
%- multilevel picard method
%- deep backward storchastic eq?
%- deep galerkin
%on the machine learning side
%
%and feynman-kac monte carlo formulation


%% This thesis goal:
This thesis investigates the use of physical-informed nural netowrks for solving
diffusion-type partial differential equations in moderate high-dimensional setting.
These problems are especially relevant
because a large subclass of problems, especially those involving 
probability density functions, often live in high-dimensional state spaces.

%% This thesis struct:
The first chapters descibe the
mathematical background of diffusion equations, stochastic differential equations,
and Fokker-Planck equations.
We follow with introducing the
sparse grid combination technique.
After that, feed-forward neural networks and physics-
informed neural networks are presented.

The numerical part of the thesis considers selected diffusion-type problems
including the heat equation, the Smoluchowski equation, and Ornstein-Uhlenbeck processes.
These results are used to evaluate the
accuracy, computational cost, scalability,
and limitations of the considered methods
in higher dimensional setting.



